\documentclass[12pt,a4paper]{ctexart}

% === 页面与字体 ===
\usepackage[top=2.5cm,bottom=2.5cm,left=2.8cm,right=2.8cm]{geometry}
\usepackage{amsmath,amssymb,amsthm}
\usepackage{mathrsfs}
\usepackage{mathtools}
\usepackage{tikz-cd}
\usepackage{hyperref}
\usepackage{enumitem}
\usepackage{titlesec}
\usepackage{fancyhdr}

% === 定理环境 ===
\newtheorem{definition}{定义}[section]
\newtheorem{theorem}{定理}[section]
\newtheorem{proposition}{命题}[section]
\newtheorem{corollary}{推论}[section]
\newtheorem{lemma}{引理}[section]
\newtheorem{remark}{注记}[section]
\newtheorem{example}{例}[section]
\theoremstyle{definition}
\newtheorem*{proof*}{证明}

% === 自定义命令 ===
\newcommand{\cat}[1]{\mathbf{#1}}
\newcommand{\Set}{\cat{Set}}
\newcommand{\Cat}{\cat{Cat}}
\newcommand{\Ob}{\mathrm{Ob}}
\newcommand{\Hom}{\mathrm{Hom}}
\newcommand{\Id}{\mathrm{Id}}
\newcommand{\id}{\mathrm{id}}
\newcommand{\op}{\mathrm{op}}
\newcommand{\nat}{\Rightarrow}
\newcommand{\liff}{\;\Longleftrightarrow\;}

% === 页眉高度 ===
\setlength{\headheight}{14pt}
\addtolength{\topmargin}{-2pt}

% === 页眉页脚 ===
\pagestyle{fancy}
\fancyhf{}
\fancyhead[L]{\small 范畴论笔记}
\fancyhead[R]{\small 严格 2-范畴中的伴随}
\fancyfoot[C]{\thepage}
\renewcommand{\headrulewidth}{0.5pt}

\title{\bfseries 严格 2-范畴中的伴随\\[0.3cm] \large —— 从代数公理到经典 Hom-集双射的等价性证明}
\author{}
\date{\today}

\begin{document}

\maketitle
\tableofcontents
\newpage

% ==================== 引言 ====================
\section{引言}

在经典范畴论中，伴随函子通过 Hom-集的自然双射来定义：
\[
\Hom_{\mathcal{D}}(FX, Y) \cong \Hom_{\mathcal{C}}(X, GY).
\]
这一概念具有极强的统一性，连接了自由-遗忘、积-指数、张量-Hom 等大量数学结构。
但伴随还有另一个等价刻画——通过单位元 $\eta: \Id_{\mathcal{C}} \Rightarrow GF$ 和上单位元 $\varepsilon: FG \Rightarrow \Id_{\mathcal{D}}$ 配合三角恒等式来定义。

本文将直接切入范畴论的严格公理体系，分三步走：
\begin{enumerate}[leftmargin=*]
    \item 给出\textbf{严格 2-范畴（Strict 2-Category）}的代数定义；
    \item 在任意 2-范畴中，纯代数地定义\textbf{伴随（Adjunction）}；
    \item 以 $\Cat$ 为特例，严格证明\textbf{2-范畴中的伴随等价于经典 Hom-集双射}。
\end{enumerate}

本笔记力求摒弃冗余比喻，直接触碰最硬核的数学结构。

% ==================== 2-范畴的严谨定义 ====================
\section{严格 2-范畴的严谨定义}

\begin{definition}[严格 2-范畴]
一个\textbf{严格 2-范畴（Strict 2-Category）} $\mathcal{K}$ 是一个包含三层结构和两种合成方式的代数结构，具体由以下数据构成。
\end{definition}

\subsection{0-细胞：对象类}

$\mathcal{K}$ 的\textbf{对象}（或称\textbf{0-细胞}）记作 $\Ob(\mathcal{K})$ 或 $\mathcal{K}_0$，其元素记为
\[
A, B, C, \dots
\]

\subsection{1-细胞与 2-细胞：态射范畴}

对于任意两个对象 $A, B \in \mathcal{K}_0$，存在一个小范畴 $\mathcal{K}(A, B)$（称为\textbf{Hom-范畴}）：
\begin{itemize}[leftmargin=*]
    \item $\mathcal{K}(A, B)$ 的\textbf{对象}称为\textbf{1-细胞}，记作 $f, g: A \to B$；
    \item $\mathcal{K}(A, B)$ 的\textbf{态射}称为\textbf{2-细胞}，记作 $\alpha: f \Rightarrow g$。
\end{itemize}

\subsection{垂直合成（Vertical Composition）}

由于 $\mathcal{K}(A, B)$ 本身是一个范畴，2-细胞之间可以进行范畴内的常规合成。
对于
\[
\alpha: f \Rightarrow g,\qquad \beta: g \Rightarrow h,
\]
存在\textbf{垂直合成} $\beta \cdot \alpha: f \Rightarrow h$。它满足结合律，且每个 1-细胞 $f$ 都有恒等 2-细胞
\[
1_f: f \Rightarrow f.
\]

\subsection{水平合成（Horizontal Composition）}

存在一个\textbf{双函子（bifunctor）}
\[
\ast: \mathcal{K}(B, C) \times \mathcal{K}(A, B) \longrightarrow \mathcal{K}(A, C),
\]
满足：
\begin{itemize}[leftmargin=*]
    \item 将 1-细胞 $f: A \to B$ 和 $g: B \to C$ 映射为 1-细胞 $g \circ f: A \to C$（通常仍记为 $gf$ 或 $g \ast f$）；
    \item 将 2-细胞 $\alpha: f_1 \Rightarrow f_2$ 和 $\beta: g_1 \Rightarrow g_2$ 映射为 2-细胞
    \[
    \beta \ast \alpha: g_1 \circ f_1 \Rightarrow g_2 \circ f_2.
    \]
\end{itemize}

\subsection{核心公理：互换律（Interchange Law）}

由于水平合成 $\ast$ 被定义为\textbf{双函子}，这自动强制要求它在两个变元上联合自然。
这给出了 2-范畴中最核心的代数等式——\textbf{互换律}。

对于任意合法的 2-细胞 $\alpha, \alpha', \beta, \beta'$，必须满足：
\begin{equation}\label{eq:interchange}
(\beta' \cdot \beta) \ast (\alpha' \cdot \alpha) = (\beta' \ast \alpha') \cdot (\beta \ast \alpha).
\end{equation}

换言之，在下面的图表中，水平合成与垂直合成的顺序可以互换：
\[
\begin{tikzcd}[column sep=large, row sep=large]
A \ar[r, bend left=60, ""{name=U,below}] \ar[r, bend right=60, ""{name=D,above}]
& B \ar[r, bend left=60, ""{name=U2,below}] \ar[r, bend right=60, ""{name=D2,above}]
& C
\arrow[Rightarrow, from=U, to=D, "\alpha"'] 
\arrow[Rightarrow, from=U2, to=D2, "\beta"']
\end{tikzcd}
\]

% ==================== 2-范畴中的伴随 ====================
\section{2-范畴中的“伴随”严格定义}

在上述严谨的 2-范畴 $\mathcal{K}$ 中，我们完全抛开集合、抛开 Hom-集、抛开元素，以纯代数的方式定义伴随。

\begin{definition}[伴随（Adjunction in a 2-Category）]
设 $A, B \in \mathcal{K}_0$ 为两个 0-细胞。一个\textbf{伴随（Adjunction）} 由四元组 $(L, R, \eta, \varepsilon)$ 构成，其中：
\begin{enumerate}[leftmargin=*]
    \item \textbf{1-细胞：} $L: A \to B$ 和 $R: B \to A$；
    \item \textbf{2-细胞：} 
    \begin{align*}
        \eta &: 1_A \Rightarrow R \ast L \quad \text{（单位元，Unit）},\\
        \varepsilon &: L \ast R \Rightarrow 1_B \quad \text{（上单位元，Counit）}.
    \end{align*}
    这里的 $1_A$ 是 $A \to A$ 的恒等 1-细胞，$1_B$ 是 $B \to B$ 的恒等 1-细胞。
\end{enumerate}
\end{definition}

\subsection{三角恒等式（Triangle Identities）}

四元组必须满足以下两个\textbf{2-细胞方程}（使用垂直合成 $\cdot$ 和水平合成 $\ast$）：

\begin{align}
(1_R \ast \varepsilon) \cdot (\eta \ast 1_R) &= 1_R, \label{eq:triangle1}\\[4pt]
(\varepsilon \ast 1_L) \cdot (1_L \ast \eta) &= 1_L. \label{eq:triangle2}
\end{align}

这两个恒等式可以用交换图表直观表示为：

\[
\begin{tikzcd}[column sep=large]
R \ar[r, "\eta \ast 1_R"] \ar[dr, "1_R"'] & RLR \ar[d, "1_R \ast \varepsilon"] \\
& R
\end{tikzcd}
\qquad
\begin{tikzcd}[column sep=large]
L \ar[r, "1_L \ast \eta"] \ar[dr, "1_L"'] & LRL \ar[d, "\varepsilon \ast 1_L"] \\
& L
\end{tikzcd}
\]

这就是 2-范畴中伴随的\textbf{全部定义}。没有底层元素，只有结构。

% ==================== 终极证明 ====================
\section{终极证明：从 2-范畴伴随到经典 Hom-双射}

现在我们把一般 2-范畴 $\mathcal{K}$ 特例化为\textbf{小范畴构成的 2-范畴 $\Cat$}。
这是范畴论中最重要的 2-范畴之一：

\begin{itemize}[leftmargin=*]
    \item \textbf{0-细胞} $A, B$ 就是普通范畴 $\mathcal{C}, \mathcal{D}$；
    \item \textbf{1-细胞} $L, R$ 就是函子 $F: \mathcal{C} \to \mathcal{D},\; G: \mathcal{D} \to \mathcal{C}$；
    \item \textbf{2-细胞} $\eta, \varepsilon$ 就是自然变换；
    \item \textbf{水平合成} $\ast$ 是函子作用于自然变换（例如 $G \ast \varepsilon$ 记为 $G\varepsilon$）；
    \item \textbf{垂直合成} $\cdot$ 是自然变换的常规合成 $\circ$。
\end{itemize}

\begin{theorem}[2-范畴伴随 $\liff$ 经典 Hom-集双射]
在 $\Cat$ 中，设 $(F, G, \eta, \varepsilon)$ 满足三角恒等式 (\ref{eq:triangle1}) 和 (\ref{eq:triangle2})，
则对任意 $X \in \mathcal{C}, Y \in \mathcal{D}$，存在自然双射
\[
\Hom_{\mathcal{D}}(FX, Y) \cong \Hom_{\mathcal{C}}(X, GY).
\]
\end{theorem}

\begin{proof*}
任取对象 $X \in \mathcal{C}, Y \in \mathcal{D}$。我们需要构造两个映射 $\Phi$ 和 $\Psi$，并证明它们互为逆映射。

\medskip
\noindent\textbf{步骤 1：构造映射}

定义
\begin{align}
\Phi &: \Hom_{\mathcal{D}}(FX, Y) \longrightarrow \Hom_{\mathcal{C}}(X, GY), \notag\\
\Phi(f) &= G(f) \circ \eta_X, \quad \forall\, f: FX \to Y. \label{eq:Phi}
\end{align}

定义
\begin{align}
\Psi &: \Hom_{\mathcal{C}}(X, GY) \longrightarrow \Hom_{\mathcal{D}}(FX, Y), \notag\\
\Psi(g) &= \varepsilon_Y \circ F(g), \quad \forall\, g: X \to GY. \label{eq:Psi}
\end{align}

\medskip
\noindent\textbf{步骤 2：证明 $\Psi(\Phi(f)) = f$（左逆）}

将 $\Phi(f)$ 代入 $\Psi$ 的定义：
\begin{align*}
\Psi(\Phi(f)) &= \Psi\big(G(f) \circ \eta_X\big) \\
&= \varepsilon_Y \circ F\big(G(f) \circ \eta_X\big).
\end{align*}

由于函子 $F$ 保持态射合成，拆开得：
\begin{equation}\label{eq:step1}
\Psi(\Phi(f)) = \varepsilon_Y \circ F(G(f)) \circ F(\eta_X).
\end{equation}

\textbf{关键一步：利用自然变换 $\varepsilon$ 的自然性。}

$\varepsilon: FG \Rightarrow \Id_{\mathcal{D}}$ 是自然变换。对于 $\mathcal{D}$ 中的态射 $f: FX \to Y$，自然性要求以下图表交换：
\[
\begin{tikzcd}
FG(FX) \ar[r, "FG(f)"] \ar[d, "\varepsilon_{FX}"'] & FG(Y) \ar[d, "\varepsilon_Y"] \\
FX \ar[r, "f"'] & Y
\end{tikzcd}
\]
即
\begin{equation}\label{eq:naturality-eps}
f \circ \varepsilon_{FX} = \varepsilon_Y \circ F(G(f)).
\end{equation}

将 (\ref{eq:naturality-eps}) 代入 (\ref{eq:step1})：
\begin{align*}
\Psi(\Phi(f)) &= \big[f \circ \varepsilon_{FX}\big] \circ F(\eta_X) \\
&= f \circ \big[\varepsilon_{FX} \circ F(\eta_X)\big].
\end{align*}

\textbf{绝杀一步：利用三角恒等式。}

回忆第二个三角恒等式 (\ref{eq:triangle2})：$(\varepsilon \ast 1_F) \cdot (1_F \ast \eta) = 1_F$。
在分量 $X$ 处，它精确等于：
\begin{equation}\label{eq:triangle2-component}
\varepsilon_{FX} \circ F(\eta_X) = \id_{FX}.
\end{equation}

代入得：
\[
\Psi(\Phi(f)) = f \circ \id_{FX} = f.
\]

\medskip
\noindent\textbf{步骤 3：证明 $\Phi(\Psi(g)) = g$（右逆）}

完全对称的推导。将 $\Psi(g)$ 代入 $\Phi$：
\begin{align*}
\Phi(\Psi(g)) &= \Phi\big(\varepsilon_Y \circ F(g)\big) \\
&= G\big(\varepsilon_Y \circ F(g)\big) \circ \eta_X \\
&= G(\varepsilon_Y) \circ G(F(g)) \circ \eta_X.
\end{align*}

\textbf{利用 $\eta$ 的自然性：} 对 $\mathcal{C}$ 中态射 $g: X \to GY$，$\eta$ 的自然性给出：
\[
\begin{tikzcd}
X \ar[r, "g"] \ar[d, "\eta_X"'] & GY \ar[d, "\eta_{GY}"] \\
GF(X) \ar[r, "GF(g)"'] & GF(GY)
\end{tikzcd}
\]
即
\begin{equation}\label{eq:naturality-eta}
G(F(g)) \circ \eta_X = \eta_{GY} \circ g.
\end{equation}

代入得：
\begin{align*}
\Phi(\Psi(g)) &= G(\varepsilon_Y) \circ \big(\eta_{GY} \circ g\big) \\
&= \big[G(\varepsilon_Y) \circ \eta_{GY}\big] \circ g.
\end{align*}

\textbf{利用三角恒等式：} 第一个三角恒等式 (\ref{eq:triangle1}) 在分量 $Y$ 处精确等于：
\begin{equation}\label{eq:triangle1-component}
G(\varepsilon_Y) \circ \eta_{GY} = \id_{GY}.
\end{equation}

代入得：
\[
\Phi(\Psi(g)) = \id_{GY} \circ g = g.
\]

\medskip
\noindent 综上，$\Phi$ 与 $\Psi$ 互为逆映射，因此 $\Hom_{\mathcal{D}}(FX, Y) \cong \Hom_{\mathcal{C}}(X, GY)$ 是双射。
\end{proof*}

% ==================== 结论 ====================
\section{结论与洞察}

通过上述严格证明，我们可以清晰地看到 Hom-集双射的本质来源：

\begin{enumerate}[leftmargin=*]
    \item \textbf{自然变换的自然性（交换图表）}完成了态射的\textbf{“平移”}——它允许我们在 $F$ 和 $G$ 之间穿梭，将 $\mathcal{D}$ 中的态射“搬运”到 $\mathcal{C}$ 中，反之亦然；
    
    \item \textbf{2-范畴的公理（三角恒等式）}确保了过去再回来的操作能够严格\textbf{“抵消”为恒等态射}——这正是 $\Psi \circ \Phi = \id$ 和 $\Phi \circ \Psi = \id$ 的代数根源。
\end{enumerate}

经典伴随（Hom-集双射定义）与 2-范畴伴随（单位元/上单位元 + 三角恒等式定义）的等价性，
并非巧合，而是同一个严谨代数结构在 $\Cat$ 中的\textbf{投影}。
掌握这一视角后，伴随函子的概念可以从普通的范畴自然推广到任意的 2-范畴（乃至双范畴），
为高阶范畴论和同伦类型论等更深层次的理论铺平道路。

% ==================== 附录 ====================
\appendix
\section{关键概念速查}

\begin{definition}[范畴 $\Cat$ 的 2-范畴结构]
$\Cat$ 是所有小范畴构成的 2-范畴，其中：
\begin{itemize}[leftmargin=*]
    \item 0-细胞：小范畴 $\mathcal{C}, \mathcal{D}, \dots$
    \item 1-细胞：函子 $F: \mathcal{C} \to \mathcal{D}$
    \item 2-细胞：自然变换 $\alpha: F \Rightarrow G$
    \item 垂直合成：自然变换的逐点合成 $\circ$
    \item 水平合成：函子对自然变换的作用（Whiskering）
\end{itemize}
\end{definition}

\begin{definition}[自然变换的自然性条件]
设 $F, G: \mathcal{C} \to \mathcal{D}$ 为函子，$\alpha: F \Rightarrow G$ 为自然变换。
则对 $\mathcal{C}$ 中任意态射 $f: X \to Y$，以下图表交换：
\[
\begin{tikzcd}
F(X) \ar[r, "F(f)"] \ar[d, "\alpha_X"'] & F(Y) \ar[d, "\alpha_Y"] \\
G(X) \ar[r, "G(f)"'] & G(Y)
\end{tikzcd}
\]
即 $G(f) \circ \alpha_X = \alpha_Y \circ F(f)$。
\end{definition}

\begin{definition}[三角恒等式（分量形式）]
在 $\Cat$ 中，三角恒等式 (\ref{eq:triangle1}) 和 (\ref{eq:triangle2}) 在分量上的写法为：
\begin{align*}
G(\varepsilon_Y) \circ \eta_{GY} &= \id_{GY}, \quad \forall\, Y \in \mathcal{D},\\
\varepsilon_{FX} \circ F(\eta_X) &= \id_{FX}, \quad \forall\, X \in \mathcal{C}.
\end{align*}
\end{definition}

\end{document}